%% kapa: exam=1 round=35 year=2024 period=1 form=A subject=economics subject_ko=경제학원론 seq=15 number=55
%% source: Q-Net 감정평가사 기출문제 (cst00309) artlSeq=5213229 file=raw/1008/5213229/1교시 문제지 원본.pdf
%% conversion: kordoc 4.24.51 -> tools/gen1.py (auto); review=proofread-sonnet; flags=none
\begin{question}{55}
사회후생함수에 관한 설명으로 옳지 않은 것은?
\begin{choices}
  \item 평등주의 경향이 강할수록 사회무차별 곡선은 원점에 대해 더 오목한 모양을 갖는다.
  \item 평등주의적 사회후생함수는 개인들의 효용수준의 차이를 반영해야 한다는 평등주의적 가치판단을 근거로 한다.
  \item 공리주의자의 사회후생함수는 사회구성원의 효용수준에 동일한 가중치를 부여한다.
  \item 롤즈(J. Rawls)의 가치판단을 반영한 사회무차별곡선은 $L$자 모양이다.
  \item 롤즈의 최소극대화 원칙(maxmin principle)은 한 사회에서 가장 가난한 사람의 생활 수준을 가능한 한 크게 개선시키는 것이 재분배정책의 최우선 과제라는 주장이다.
\end{choices}
\end{question}
\answer{1}
